\subsection{Hardware realisation}

The hardware platform is an ESP32-WROOM module connected to an L298N H-bridge.
The L298N introduces an approximate voltage drop of \(\SI{4}{V}\), so with a
\(\SI{12}{V}\) supply the motor sees about \(\SI{8}{V}\). The motor is equipped
with a quadrature encoder. The sampling time is \(T=\SI{0.08}{s}\).

The controller is implemented by automatic code generation from Simulink using
the Embedded Real-Time (ERT) target and the ESP32 target support. External Mode
is used for Monitor and Tune operation during experiments.

\begin{scripttable}
  \centering
  \caption{ESP32 pin assignment for the DC motor test rig.}
  \label{tab:esp32-pin-assignment}
  \begin{tabular}{ll}
    \toprule
    Signal & ESP32 pin \\
    \midrule
    PWM & GPIO 14 \\
    IN1 & GPIO 26 \\
    IN2 & GPIO 27 \\
    Encoder A & GPIO 32 \\
    Encoder B & GPIO 33 \\
    \bottomrule
  \end{tabular}
\end{scripttable}

The encoder resolution was determined by a one-revolution test as
\(44\) counts per revolution, corresponding to \(11\) encoder impulses with
x4 quadrature edge decoding. The speed conversion
\[
  \omega = \frac{2\pi \Delta N}{T\,\mathrm{CPR}}
\]
is therefore calibrated for \(\mathrm{CPR}=44\), and the hardware speed values
are reported in physical units. The measured full-duty speed of approximately
\(\SI{780}{rad/s}\) corresponds to about \(\SI{7450}{rpm}\). This is consistent
with the gearbox-free 25GA371 speed range of roughly \(\SIrange{3000}{8000}{rpm}\)
and independently supports the encoder scaling. Hardware settling times are
not stated here.
